% Informe del Proyecto 2 (entrega en PDF, enunciado §11.2). Cada sección dice quién la escribe (TODO.md).
% Compilar desde esta carpeta, dos veces para las referencias (sin instalar LaTeX):
%   docker run --rm --user "$(id -u):$(id -g)" -e HOME=/tmp -v "$PWD":/w -w /w texlive/texlive:latest-small pdflatex informe.tex
% \falta{...} marca lo que todavía no está; antes de entregar no debe quedar ninguno.
\documentclass[10pt,letterpaper]{article}

\usepackage[T1]{fontenc}
\usepackage[utf8]{inputenc}
\usepackage[spanish,es-noquoting,es-noshorthands,es-tabla]{babel}
\usepackage{lmodern}
\usepackage[margin=1.9cm,top=1.5cm,bottom=1.6cm]{geometry}
\usepackage{booktabs}
\usepackage{array}
\usepackage{xcolor}
\usepackage{graphicx}
\usepackage{fancyvrb}

\definecolor{acc}{HTML}{1F4E79}
\definecolor{gris}{HTML}{5B6672}

\DeclareUnicodeCharacter{2192}{$\rightarrow$}
\DeclareUnicodeCharacter{2264}{$\le$}

% Marca de algo pendiente (sale en rojo). Antes de entregar no debe quedar ningún uso.
\newcommand{\falta}[1]{{\color{red}\bfseries [FALTA: #1]}}
% Archivo de evidence/ (toda cifra debe poder rastrearse hasta uno).
\newcommand{\ev}[1]{{\footnotesize\texttt{evidence/\allowbreak #1}}}

\setlength{\parskip}{3pt}
\setlength{\parindent}{0pt}
\renewcommand{\arraystretch}{1.05}
\setlength{\emergencystretch}{4em}

% Secciones y pie de pagina sin paquetes externos: solo LaTeX base.
\makeatletter
\renewcommand\section{\@startsection{section}{1}{\z@}%
  {9pt \@plus 2pt \@minus 1pt}{3pt}%
  {\normalfont\large\bfseries\color{acc}}}
\renewcommand\subsection{\@startsection{subsection}{2}{\z@}%
  {6pt \@plus 1pt}{2pt}%
  {\normalfont\normalsize\bfseries\color{acc}}}
\renewcommand\@seccntformat[1]{\csname the#1\endcsname.\quad}
\renewcommand{\ps@plain}{%
  \renewcommand{\@oddhead}{}\renewcommand{\@evenhead}{}%
  \renewcommand{\@oddfoot}{\footnotesize\color{gris}%
    TI-4601 $\cdot$ Proyecto 2 $\cdot$ Plataforma distribuida de anotación de datos
    \hfill \thepage}%
  \renewcommand{\@evenfoot}{\@oddfoot}}
\makeatother
\pagestyle{plain}

\fvset{fontsize=\footnotesize,frame=leftline,framerule=0.6pt,rulecolor=\color{gris},xleftmargin=6pt}

\begin{document}

{\LARGE\bfseries\color{acc} Proyecto 2 --- Plataforma distribuida de anotación de datos}\\[2pt]
{\footnotesize\color{gris} TI-4601 Bases de Datos Avanzadas $\cdot$ Prof. Manuel Zumbado Corrales $\cdot$ Grupo 3 $\cdot$ Fabricio Mena $\cdot$ Alisson Redondo $\cdot$ Christopher Rodriguez}\\
{\footnotesize\color{gris} Repositorio: \texttt{https://github.com/ChrisRC7/Proyecto-2-Bases-de-Datos-Avanzadas}}

\vspace{4pt}\hrule\vspace{6pt}

\textbf{Entorno de todas las mediciones.} Laptop AMD Ryzen 7 5700U, 16 GB de RAM (4 CPU y 9,7 GiB para
WSL 2), Windows 11 Home con Ubuntu 24.04 en WSL 2 y Docker Desktop 29.5.2. Todos los servicios, los tres
nodos de MongoDB y el simulador corren como contenedores en esa misma máquina. Cada cifra del informe cita
el archivo de \texttt{evidence/} que la respalda; el Apéndice~A los reúne.

\section{Arquitectura y recorrido de una tarea LiDAR}

\falta{B6.1 (Alisson): arquitectura general, diagrama, límites de responsabilidad de cada servicio,
recorrido completo de una tarea LiDAR (siguiente, borrador, enviar, revisión, aprobada, visible en Q1) y
evidencia del aislamiento modular (redes, variables de entorno: el registro no tiene la dirección de
Valkey, el simulador solo está en la red \texttt{borde})}

\section{Registro: modelo, índices y asignación exclusiva}

El registro es el único servicio que habla con MongoDB y la única fuente de verdad del estado de una
tarea. No tiene la dirección de ningún Valkey: lo que necesita del lock se lo pasa el gateway (el token).

\subsection{Modelo}

MongoDB corre como replica set \texttt{rs0} de tres nodos (\texttt{mongo:7.0}) y el registro escribe con
\texttt{w: "majority"}. La base \texttt{plataforma} tiene las colecciones \texttt{clientes},
\texttt{proyectos}, \texttt{batches}, \texttt{anotadores} y \texttt{tareas}, todas con UUID como
\texttt{\_id}. El documento de una tarea guarda la referencia al archivo (\texttt{uri}, \texttt{sha256},
\texttt{bytes}, \texttt{mime}, \texttt{meta}), nunca el archivo; además el \texttt{estado}, la
\texttt{version}, la \texttt{asignacion} vigente (anotador, equipo, token, \texttt{tomada\_en},
\texttt{vence\_en}), la \texttt{anotacion} vigente, la \texttt{ultima\_revision} y el
\texttt{ultimo\_evento\_id}. El \texttt{cliente\_id} y el \texttt{sla\_horas} del proyecto se copian en la
tarea para armar el evento del lago sin otra consulta.

Cada transición es un solo \texttt{findOneAndUpdate} cuyo filtro lleva el estado requerido (en
\texttt{enviar} y \texttt{liberar}, también el token de la asignación; y la versión si viene en la petición). Si el filtro no encuentra el documento,
la respuesta es 409 y no se escribe nada, tampoco en el lago. Si lo encuentra, cambia el estado, suma 1 a
\texttt{version} y genera el \texttt{event\_id} que el gateway manda al historial.

\begin{center}\footnotesize
\begin{tabular}{@{}lllll@{}}
\toprule
Ruta del gateway & Acción (registro) & Estado requerido & MongoDB queda & Fila del lago \\
\midrule
\texttt{POST /tareas/siguiente} & asignar & \texttt{disponible} o \texttt{abierta} vencida & \texttt{abierta} & \texttt{abierta} \\
\texttt{POST /tareas/\{id\}/enviar} & enviar & \texttt{abierta} + token & \texttt{anotada} & \texttt{anotada} \\
\texttt{POST /tareas/\{id\}/revision} & empezar\_revision & \texttt{anotada} & \texttt{en\_revision} & \texttt{en\_revision} \\
\texttt{POST /tareas/\{id\}/revision} & aprobar & \texttt{en\_revision} & \texttt{aprobada} & \texttt{aprobada} \\
\texttt{POST /tareas/\{id\}/revision} & rechazar (con motivo) & \texttt{en\_revision} & \texttt{disponible} & \texttt{rechazada} \\
\texttt{POST /tareas/\{id\}/liberar} & liberar & \texttt{abierta} + token & \texttt{disponible} & \texttt{liberada} \\
\bottomrule
\end{tabular}
\end{center}

La última fila no está en la tabla del enunciado: la agregamos porque la ruta \texttt{liberar} sí cambia el
estado, y así queda trazada. Las consultas Q1--Q5 no la usan. En el rechazo, MongoDB devuelve la tarea a
\texttt{disponible} y borra la asignación, y la fila del lago queda \texttt{rechazada} con su código de
motivo (uno de cinco: \texttt{ETIQUETA\_INCORRECTA}, \texttt{GEOMETRIA\_IMPRECISA},
\texttt{OBJETO\_OMITIDO}, \texttt{INSTRUCCION\_NO\_SEGUIDA}, \texttt{OTRO}); el gateway responde 422 si
falta. La prueba de las rutas recorre una tarea LiDAR aprobada y otra rechazada, con todos los códigos de
error del contrato (\ev{A4\_rutas\_gateway\_20261110-1446.txt} y
\ev{A\_correcciones\_gateway\_20261110-1713.txt}).

\textbf{Datos.} \texttt{dataset/generar.py} crea con la semilla 2026 los volúmenes del enunciado: 60
clientes, 120 proyectos activos (los 60 cerrados solo existen en el historial), 600 batches, 60\,000 tareas
(30\,\% texto, 30\,\% imagen, 20\,\% video y 20\,\% LiDAR, cada una con su tipo de geometría) y 300
anotadores en 30 equipos, 30 de ellos revisores. El primer proyecto es el del enunciado: LiDAR urbano 2026,
SLA de 48 h y la taxonomía \texttt{lidar\#peaton}, \texttt{lidar\#vehiculo}, \texttt{lidar\#ciclista} (un
skater es ciclista) y \texttt{lidar\#signal}. \texttt{docker compose --profile datos run --rm dataset}
borra y vuelve a cargar todo en unos 5~s, con los mismos ids en cualquier máquina (\texttt{datos/README.md}).

\subsection{Índices}

Se crean los cuatro obligatorios: en \texttt{tareas}, \texttt{\{proyecto\_id, estado, prioridad\}},
\texttt{\{asignacion.anotador\_id, estado\}} y \texttt{\{batch\_id\}}; en \texttt{proyectos},
\texttt{\{cliente\_id, estado\}}. Para decidir si hacía falta alguno más se corrió
\texttt{explain("executionStats")} de las dos consultas que más se repiten, la asignación y la exportación
(\ev{A3\_explain\_registro\_20261022-2003.md}):

\begin{center}\small
\begin{tabular}{@{}llrrr@{}}
\toprule
Consulta & Índices & Claves & Documentos & Devueltos \\
\midrule
Asignación con \texttt{proyecto\_id} & obligatorios & 1 & 1 & 1 \\
Asignación sin \texttt{proyecto\_id} (cola global) & obligatorios & 0 & 60\,000 & 1 \\
Asignación sin \texttt{proyecto\_id} (cola global) & + \texttt{\{estado, prioridad\}} & 1 & 1 & 1 \\
Exportación, primera página de 100 & obligatorios & 300 & 300 & 100 \\
\bottomrule
\end{tabular}
\end{center}

Con proyecto, el índice obligatorio basta: cada rama del \texttt{\$or} (\texttt{disponible}, o
\texttt{abierta} vencida) es un rango del índice y \texttt{SORT\_MERGE} las une ya ordenadas por prioridad.
Sin proyecto, que el contrato permite y que es lo que usa el simulador, ningún índice obligatorio empieza
por \texttt{estado} y cada asignación recorría las 60\,000 tareas (\texttt{COLLSCAN} + \texttt{SORT}). Por
eso se agregó \texttt{tareas \{estado: 1, prioridad: -1\}}, que lo deja en un documento examinado. La
exportación no lleva índice propio: ordena en memoria las aprobadas de un proyecto, que son a lo sumo 500, y
un índice más se actualizaría en cada transición, porque todas cambian \texttt{estado}.

\subsection{Asignación exclusiva: 50 solicitudes contra 10 tareas}

La asignación es una sola operación sobre un documento:

\begin{Verbatim}
tareas.find_one_and_update(
    {"$or": [{"estado": "disponible"}, {"estado": "abierta", "asignacion.vence_en": {"$lt": ahora}}]},
    {"$set": {"estado": "abierta", "asignacion": {..., "token": token}}, "$inc": {"version": 1}},
    sort=[("prioridad", -1)])
\end{Verbatim}

MongoDB aplica cada actualización de un documento de forma atómica. Si dos peticiones apuntan a la misma
tarea, solo una la encuentra en \texttt{disponible}; la otra sigue con la próxima tarea o recibe 204.
El gateway usa el mismo token para el lock de Valkey y, si el lock falla, devuelve la tarea con
\texttt{liberar}. Pero la exclusividad no depende del lock, sino de este filtro.

La prueba P2 crea por el gateway un proyecto con un solo lote de exactamente 10 tareas y lanza 50
solicitudes desde 50 procesos (\texttt{ProcessPoolExecutor}) que cruzan una \texttt{Barrier} juntos; salieron
en una ventana de 68~ms. Corre en el contenedor del simulador, que solo ve al gateway
(\ev{P2\_concurrencia\_gateway\_20261110-1714.md}):

\begin{center}\small
\begin{tabular}{@{}lr@{}}
\toprule
Resultado & Valor \\
\midrule
Respuestas 200 (tarea asignada) / 204 (sin tarea) & 10 / 40 \\
Tareas con más de un anotador & 0 \\
Envío con un token ajeno $\rightarrow$ 409 & 10 de 10 \\
Primer envío $\rightarrow$ 200; segundo envío de la misma tarea $\rightarrow$ 409 & 10 de 10; 10 de 10 \\
Exportación: aprobadas con el anotador que las tomó & 10 de 10 \\
Latencia en el cliente, n = 50: p50 / p95 / p99 & 324 / 500 / 510 ms \\
\texttt{POST /tareas/asignar} en el registro, n = 50: p50 / p95 / p99 & 16 / 82 / 93 ms \\
\bottomrule
\end{tabular}
\end{center}

El token ajeno lo rechaza la sesión (no coincide con el lock) antes de llegar a MongoDB. El segundo envío lo
rechaza el registro porque la tarea ya no está en \texttt{abierta}, y como no hubo transición tampoco se
escribe fila en el lago. La misma prueba directa contra el registro, sin gateway, dio el mismo resultado
(\ev{P2\_concurrencia\_registro\_20261022-2007.md}). La latencia se explica en la sección~6: antes de
corregir el gateway, esta misma prueba tenía p50 de 3,4~s en el cliente.

\subsection{\texttt{w: "majority"} frente a \texttt{w: 1}}

Se midió con la misma carga y el mismo dataset, cambiando solo \texttt{MONGO\_W} en el registro y
llamándolo directo, sin gateway ni Valkey, para aislar MongoDB. Cada ciclo hace cuatro escrituras (asignar,
enviar, empezar revisión, aprobar) y el dataset se recargó antes de cada corrida
(\ev{A7.4\_write\_concern\_20261022-2011.md}):

\begin{center}\small
\begin{tabular}{@{}lrr@{}}
\toprule
 & \texttt{majority} & \texttt{w: 1} \\
\midrule
1 hilo: \texttt{enviar} p50 / p95 / p99 & 11,8 / 17,1 / 23,5 ms & 8,8 / 13,3 / 20,3 ms \\
1 hilo: ciclos por segundo & 20,9 & 27,6 \\
20 hilos (3 corridas de 30 s): \texttt{enviar} p50 / p95 / p99 & 133,0 / 182,2 / 212,1 ms & 129,2 / 171,9 / 196,6 ms \\
20 hilos: ciclos por segundo (corridas 1, 2, 3) & 39,0 $\cdot$ 37,9 $\cdot$ 38,4 & 39,9 $\cdot$ 39,8 $\cdot$ 39,3 \\
\bottomrule
\end{tabular}
\end{center}

Sin cola (1 hilo) se ve el costo real: con \texttt{majority} el primario espera a que un secundario aplique
la entrada del oplog, unos 3~ms más por escritura y 24\,\% menos ciclos por segundo. Con 20 hilos la
diferencia baja a un 3\,\% porque manda la cola: el registro es un solo proceso de Python y el primario
también usa casi un núcleo, con 4 CPU para todo. Esto no se puede llevar a un clúster real. Aquí los tres
nodos están en la misma máquina y la réplica viaja por la red virtual de Docker; con nodos en otras máquinas,
\texttt{majority} sumaría al menos un viaje de red por escritura y la diferencia sería mayor. Lo que
\texttt{w: 1} ahorra en latencia lo pierde en durabilidad: una escritura confirmada solo por el primario se
puede perder si ese primario cae antes de replicarla (sección~7.1).

\subsection{Límites del registro}

\begin{itemize}\setlength{\itemsep}{1pt}
\item El plazo de la asignación (\texttt{vence\_en}) es fijo: 120~s desde que se toma. El lock de Valkey se
renueva con cada borrador y latido, pero el plazo en MongoDB no. Si alguien trabaja más de 120~s, el
registro puede ofrecer la tarea a otra persona mientras el lock sigue tomado; el gateway no logra el lock,
devuelve la tarea y responde 503. Lo anotamos como tarea pendiente; en las pruebas nadie tarda tanto.
\item La exportación ordena en memoria; está bien para 500 aprobadas por proyecto, no para miles.
\item Los explain y las mediciones son de una sola máquina con una colección pequeña: comparan documentos
examinados y latencias relativas, no tiempos de un despliegue real.
\end{itemize}


\section{Sesión, borrador y reinicio de Valkey}

\falta{B6.2 (Alisson): llaves de Valkey, lock con SET NX PX, borrador cada 2 s, recuperación de la
sesión y qué pasa al reiniciar Valkey con AOF everysec (P3)}

\section{Lago Parquet}

\falta{C6.1 (Christopher): micro-lotes con \_tmp y rename, compactación, duplicados por event\_id, la
fila \texttt{rechazada} en Parquet mientras MongoDB queda en \texttt{disponible}, frescura (P4)}

\section{Analítica con Spark}

\falta{C6.2 (Christopher): definición de Q1--Q5, resultados validados, planes de Q1 y Q3, primera
ejecución contra las siguientes, uno y dos workers, principal y grande, limitaciones de Spark (P5)}

\section{Carga y cuellos de botella}

\subsection{Cómo se midió}

\texttt{pruebas/p6\_carga.sh} corre el simulador con 10, 40 y 100 usuarios virtuales: tres corridas por
nivel, cada una con 30~s de calentamiento y 90~s medidos. Antes de cada corrida recarga el dataset y borra
de Valkey los locks, sesiones y borradores de la corrida anterior, para que las nueve empiecen igual.
Mientras corre, guarda \texttt{docker stats} cada $\sim$2~s.

El simulador corre en su contenedor, en la red \texttt{borde}, y solo habla con el gateway. Cada usuario es
un hilo. El 90\,\% son anotadores: piden tarea de la cola global, tardan entre 0,5 y 2~s en anotar y
envían. El 10\,\% son revisores: toman una tarea enviada, empiezan la revisión y la aprueban, o la rechazan
con un motivo al azar el 20\,\% de las veces. El cliente espera cada respuesta hasta 10~s. Cada petición
queda en un CSV con su fase, endpoint, código y milisegundos. Los percentiles son por rango más cercano,
juntando las tres corridas de cada nivel (solo la fase medida). P6 se corrió antes de integrar los
borradores de B4.2 (opción \texttt{--borradores}), así que esta carga no incluye
\texttt{PUT /tareas/\{id\}/borrador} ni los latidos.

\subsection{Resultados}

Corrida final, con el gateway corregido (\ev{P6\_carga\_20261110-1800.md}):

\begin{center}\small
\begin{tabular}{@{}rlrrrrr@{}}
\toprule
Usuarios & Endpoint & n & p50 (ms) & p95 (ms) & p99 (ms) & Sin respuesta \\
\midrule
10 & \texttt{POST /tareas/siguiente} & 1\,660 & 61,7 & 167,3 & 250,4 & 0 \\
10 & \texttt{POST /tareas/\{id\}/enviar} & 1\,686 & 117,5 & 297,3 & 395,6 & 0 \\
10 & \texttt{POST /tareas/\{id\}/revision} & 765 & 41,9 & 110,2 & 162,1 & 0 \\
\midrule
40 & \texttt{POST /tareas/siguiente} & 2\,910 & 681,4 & 1\,448,0 & 2\,005,9 & 0 \\
40 & \texttt{POST /tareas/\{id\}/enviar} & 2\,990 & 1\,297,7 & 2\,406,1 & 2\,911,6 & 0 \\
40 & \texttt{POST /tareas/\{id\}/revision} & 1\,812 & 273,9 & 533,5 & 775,5 & 0 \\
\midrule
100 & \texttt{POST /tareas/siguiente} & 2\,526 & 2\,868,9 & 6\,046,6 & 8\,393,2 & 5 \\
100 & \texttt{POST /tareas/\{id\}/enviar} & 2\,705 & 5\,177,1 & 9\,383,6 & 10\,011,8 & 93 \\
100 & \texttt{POST /tareas/\{id\}/revision} & 2\,542 & 702,2 & 1\,452,3 & 2\,055,5 & 0 \\
\bottomrule
\end{tabular}

\vspace{4pt}
\begin{tabular}{@{}rrrrrr@{}}
\toprule
Usuarios & Envíos/min (3 corridas) & Media & Peticiones/s & CPU del gateway, media / máx. & Atraso al lago, p50 \\
\midrule
10 & 381 $\cdot$ 369 $\cdot$ 374 & 374,7 & 15,2 & 40 / 64\,\% & 2,6--2,7 s \\
40 & 667 $\cdot$ 655 $\cdot$ 672 & 664,4 & 28,6 & 93 / 114\,\% & 3,0--3,1 s \\
100 & 587 $\cdot$ 592 $\cdot$ 563 & 580,4 & 28,8 & 99 / 123\,\% & 3,6--3,8 s \\
\bottomrule
\end{tabular}
\end{center}

No hubo ningún 4xx ni 5xx en las nueve corridas. Las 98 peticiones «sin respuesta» con 100 usuarios son
las que el simulador dejó de esperar a los 10~s. El gateway también las terminó con 200: en el lago, los
eventos distintos de cada corrida son exactamente las respuestas 200 más esas peticiones (por ejemplo
3\,365 + 37 = 3\,402), sin perder ni duplicar ninguno (\texttt{lago.txt}). El CPU está en \% de un núcleo,
y la máquina tiene 4. El atraso al lago es el tiempo desde la transición en el registro hasta que el
archivo Parquet queda publicado; con micro-lotes de 5~s, unos 2,6~s es el mínimo.

\subsection{Cuellos de botella}

\textbf{Encontrados y corregidos antes de medir.} La primera vez que el simulador pasó por el gateway, 79
de 139 pedidos de tarea dieron 503 con solo 10 usuarios (\ev{A6.2\_simulador\_gateway\_20261110-1501.txt}).
Con 10 usuarios y el mismo comando quedó en 0 errores y 372 envíos por minuto
(\ev{A6.2\_simulador\_gateway\_20261110-1726.txt}) después de corregir, en la rama de integración, cinco
cosas del gateway (\ev{A\_correcciones\_gateway\_20261110-1713.txt}):
\begin{itemize}\setlength{\itemsep}{1pt}
\item \textbf{Un cliente HTTP nuevo en cada llamada interna.} Crear uno cuesta de 9 a 35~ms de CPU según la
carga, y ese trabajo bloquea el único proceso del gateway. En la P2 el registro respondía en 52~ms y el
usuario esperaba 3,4~s; con un cliente compartido, 0,32~s.
\item \textbf{Limpiezas en segundo plano.} El lock de una tarea enviada se soltaba cada 2~s desde un
trabajador. Si el revisor la rechazaba antes, la tarea volvía a la cola con el lock viejo, y cada
\texttt{/tareas/siguiente} la tomaba, fallaba el lock y respondía 503. Ahora se suelta antes de responder.
\item \textbf{Limpiezas reintentadas para siempre.} Las que recibían 404 o 409, que ya no se pueden completar,
se reprogramaban sin fin: había 637 atascadas. Ahora se cierran.
\item \textbf{SQLite bloqueado.} Con varias escrituras a la vez aparecían esperas de 10~s y
\texttt{database is locked}. Con WAL cada commit cuesta $\sim$10~ms en lugar de $\sim$22~ms, leer no bloquea a
quien escribe y las conexiones se cierran siempre (\ev{A6.2\_simulador\_gateway\_20261110-1721.txt}).
\item \textbf{Regla de \texttt{/salud}:} no afectaba la carga, pero marcaba \texttt{degradado} con MongoDB
sin primario; ahora es \texttt{caido}.
\end{itemize}

\textbf{Encontrado con la primera P6: el publicador de eventos.} En la primera corrida completa
(\ev{P6\_carga\_20261110-1727.md}) los eventos se mandaban al historial de a uno por POST. Con 100 usuarios
se generaban 27--29 eventos por segundo y llegaban 8--11, y el atraso al lago fue de 3 a 6 minutos (p50 de
199 a 324~s). Ningún evento se perdía, pero Q1 los habría visto muy tarde. Ahora van hasta 50 en un solo
POST, y el atraso con 100 usuarios es de 3,7~s. La corrida final da 6--12\,\% menos envíos por minuto que la
primera con 40 y 100 usuarios: antes el gateway dejaba ese trabajo para después de la medición, y ahora lo
hace en el momento.

\textbf{El que queda: el proceso del gateway.} De 10 a 40 usuarios las peticiones por segundo casi se
duplican (15,2 $\rightarrow$ 28,6); de 40 a 100 no suben (28,8). El gateway está en 93--99\,\% de un núcleo
y cada petición pública le cuesta unos 26--34~ms de CPU: sus llamadas internas, guardar en SQLite y mandar el
evento. Con un solo proceso de uvicorn no pasa de $\sim$30 peticiones por segundo. MongoDB y Valkey no son el
límite: con 100 usuarios el registro resuelve la transición en p50 40~ms y la sesión el lock en 25~ms,
mientras el cliente espera 5,1~s el envío (\texttt{desglose.txt}). Las cinco llamadas internas del envío
suman $\sim$150~ms; el resto es esperar turno dentro del gateway. Los tres nodos de MongoDB suman
$\sim$100\,\% de CPU y ninguno se satura. El envío es la ruta más cara porque hace más trabajo dentro del
gateway: renovar el lock, la transición, tres escrituras en SQLite con \texttt{fsync} y tres limpiezas en
sesión.

Para comprobarlo se repitió una corrida de 40 y una de 100 usuarios con el gateway en 2 procesos
(\texttt{uvicorn --workers 2}, sin cambiar código; \ev{P6\_gateway\_2procesos\_20261110-1820/}):

\begin{center}\small
\begin{tabular}{@{}lrrrr@{}}
\toprule
 & 40 u, 1 proceso & 40 u, 2 procesos & 100 u, 1 proceso & 100 u, 2 procesos \\
\midrule
Envíos/min & 664 & 917 & 580 & 887 \\
\texttt{enviar} p50 / p99 (ms) & 1\,298 / 2\,912 & 637 / 3\,211 & 5\,177 / 10\,012 & 3\,242 / 7\,960 \\
Sin respuesta en 10 s & 0 & 0 & 98 & 2 \\
\bottomrule
\end{tabular}
\end{center}

Con dos procesos salen 38--53\,\% más envíos por minuto. No llega al doble porque con 100 usuarios la
máquina completa queda en $\sim$320\,\% de sus 400\,\%: el siguiente límite sería el CPU de una sola máquina
compartido por los tres nodos de MongoDB, los servicios y el simulador. No lo dejamos como configuración
porque cada proceso correría su propio publicador de eventos y su trabajador de limpiezas sobre el mismo
SQLite, y eso solo se probó en este experimento.

\subsection{Límites de la medición}

\begin{itemize}\setlength{\itemsep}{1pt}
\item Todo corre en una laptop: el simulador compite por CPU con los servicios, y la red entre contenedores
es virtual. Los números sirven para comparar niveles y encontrar el cuello de botella, no para estimar un
despliegue en varias máquinas.
\item P6 no usa los borradores cada 2~s de B4.2; con ellos la carga de un anotador sería mayor en sesión y
en el gateway. Una corrida corta con borradores (10 usuarios, 30~s) dio 113 borradores, 113 latidos y 99
revisiones, todas con 200 (\ev{A6.2\_simulador\_borradores\_20261114-2318.txt}).
\item Tres corridas de 90~s por nivel (los envíos por minuto varían menos de $\pm$4\,\% entre ellas) y una
sola corrida en el experimento de 2 procesos.
\item Con 10\,\% de revisores se revisa alrededor de la mitad de lo que se envía; las tareas no se agotaron en
ninguna corrida (0 respuestas 204).
\end{itemize}


\section{Fallas, recuperación, RTO y RPO}

\subsection{MongoDB: cae un secundario, el primario y dos nodos}

\textbf{Configuración.} Replica set de tres nodos con voto, \texttt{electionTimeoutMillis} de 10~s (por
defecto). El registro escribe con \texttt{w: "majority"} y \texttt{wtimeout} de 5~s, y deja de esperar un
primario a los 2~s (\texttt{serverSelectionTimeoutMS}); el gateway corta cada llamada interna a 1~s. Sin
primario el registro responde 503 (\texttt{mongodb\_no\_disponible}); si la mayoría no confirma, 503
(\texttt{mongodb\_sin\_mayoria}). El gateway convierte los dos en \texttt{503 servicio\_no\_disponible}.

\textbf{Método.} \texttt{pruebas/p8\_mongo.sh} recarga el dataset y arranca una sonda en el contenedor
del simulador que, como un usuario, pide una tarea y la envía por el gateway cada $\sim$0,2~s y anota cada
respuesta en un CSV. A los 8~s mata los nodos con \texttt{docker compose kill} (SIGKILL) y guarda la marca de
tiempo. A los 15~s guarda el estado del conjunto, vuelve a encender los nodos y espera 1 PRIMARY + 2
SECONDARY. Después \texttt{p8\_verificar.py} busca en MongoDB cada envío que respondió 200. Definimos:
\begin{itemize}\setlength{\itemsep}{1pt}
\item \textbf{RPO}: envíos que respondieron 200 y no están en MongoDB, o no están en la versión esperada
(si estuvieran aplicados dos veces, la versión sería otra).
\item \textbf{RTO}: desde la marca de la falla hasta el primer envío con 200 después del último error,
tomados de filas del CSV.
\end{itemize}

\textbf{Resultado} (\ev{P8\_mongo\_gateway\_20261110-1715.md}):

\begin{center}\footnotesize
\begin{tabular}{@{}>{\raggedright}p{0.25\textwidth}>{\raggedright}p{0.2\textwidth}>{\raggedright}p{0.22\textwidth}>{\raggedright\arraybackslash}p{0.25\textwidth}@{}}
\toprule
 & Cae un secundario & Cae el primario & Caen dos (sin mayoría) \\
\midrule
Durante la falla & sigue el mismo PRIMARY & otro PRIMARY, term 14 $\rightarrow$ 15 & el que queda pasa a SECONDARY \\
Errores después de la marca & 0 & 8 $\times$ 503 & 27 $\times$ 503 \\
Envíos confirmados (200) & 122 & 132 & 164 \\
Confirmados en MongoDB, versión esperada & 122 de 122 & 132 de 132 & 164 de 164 \\
\textbf{RPO} & \textbf{0} & \textbf{0} & \textbf{0} \\
Errores que sí quedaron escritos & --- & 0 & 10 asignaciones \\
\textbf{RTO} & no se interrumpe & \textbf{10,45 s} & 32,40 s; \textbf{10,24 s} desde que vuelven los nodos \\
\bottomrule
\end{tabular}
\end{center}

\begin{itemize}\setlength{\itemsep}{1pt}
\item \textbf{Secundario.} Quedan 2 de 3 nodos, que siguen siendo mayoría: \texttt{majority} se cumple con el
primario y el otro secundario, y la escritura no se interrumpe.
\item \textbf{Primario.} Los secundarios dejan de recibir sus latidos, esperan los 10~s de
\texttt{electionTimeoutMillis} y uno gana la elección con 2 de 3 votos. El RTO (10,45~s) es casi exactamente
ese plazo: lo manda la elección, no el gateway. Directo al registro dio 10,36~s
(\ev{P8\_mongo\_20261022-2028.md}). Lo que respondió 200 sigue en MongoDB una sola vez después de la
elección.
\item \textbf{Sin mayoría.} El primario que queda solo se degrada a SECONDARY y nada se confirma: el sistema
prefiere no escribir antes que escribir sin réplica. El RTO total depende de cuándo la prueba enciende los
nodos (22~s después de la falla). Lo que mide al sistema son los 10,24~s desde que se encienden hasta el
primer 200: arrancan, se ponen al día y hay otra elección (term 16).
\item \textbf{Error no es lo mismo que no escrito.} Antes de degradarse, el primario aplica la asignación en
su copia y espera la mayoría hasta el \texttt{wtimeout}. El gateway corta al segundo, la sonda vuelve a
pedir y cada intento deja otra tarea \texttt{abierta}: 10 en esta corrida. Al volver la mayoría ese nodo
siguió de primario y las conservó. No se pierden: su plazo vence a los 120~s y vuelven a ofrecerse. Por eso
el RPO se cuenta solo sobre lo que respondió 200. Con \texttt{w: 1} el riesgo sería el contrario (no lo
medimos): un 200 que solo estaba en el primario caído podría perderse en la elección (\emph{rollback}).
\end{itemize}

\textbf{Recuperación.} Cada escenario termina con \texttt{docker compose start} de los nodos y espera a 1
PRIMARY + 2 SECONDARY (sección «Después» de cada \texttt{.txt}); el siguiente empieza con el dataset
recargado. \texttt{pruebas/estado\_rs.sh} muestra el estado del conjunto en cualquier momento.

\textbf{Límites.} Los tres nodos están en la misma máquina y la falla es \texttt{docker kill}, no una
partición de red ni un disco lento. Hay una sola sonda, y el gateway no le dice al usuario cuál de los dos
503 del registro fue; eso solo queda en el log del registro.

\subsection{Valkey y sesión}

\falta{B6.3 (Alisson): reiniciar Valkey durante carga, interrumpir una sesión, RTO y RPO de los
borradores con AOF everysec (P8 de Valkey)}


\section{Limitaciones de la solución}

\falta{C6.3 (Christopher): consolidar las limitaciones de todos. Las del registro, la carga y las
fallas de MongoDB están al final de las secciones 2, 6 y 7.1}

\appendix
\section{Evidencia citada}

Cada archivo está en \texttt{evidence/} del repositorio. Los \texttt{.md} siguen la plantilla
(\texttt{evidence/\_PLANTILLA.md}): fecha y hora, objetivo, configuración, comando, resultado,
interpretación y archivos crudos. Los \texttt{.txt} y \texttt{.csv} son la salida original.

\begin{center}\small
\begin{tabular}{@{}p{0.47\textwidth}p{0.47\textwidth}@{}}
\toprule
Archivo & Qué respalda \\
\midrule
\texttt{A3\_explain\_registro\_20261022-2003.md} (y \texttt{.txt}) & Índices: explain de la asignación y la exportación (sección 2.2). \\
\texttt{A4\_rutas\_gateway\_20261110-1446.txt} & Rutas de proyectos, lotes, exportación y revisión por el gateway, con sus errores (2.1). \\
\texttt{P2\_concurrencia\_gateway\_20261110-1714.md} (\texttt{.txt}, \texttt{.csv}) & P2: 50 solicitudes contra 10 tareas, cero dobles, 409 (2.3). \\
\texttt{P2\_concurrencia\_gateway\_20261110-1450.md} & La misma P2 antes de corregir el gateway (p50 3,4~s) (2.3 y 6.3). \\
\texttt{P2\_concurrencia\_registro\_20261022-2007.md} & P2 directa contra el registro, sin gateway (2.3). \\
\texttt{A7.4\_write\_concern\_20261022-2011.md} (\texttt{.txt}, \texttt{.csv}) & \texttt{w: "majority"} frente a \texttt{w: 1} (2.4). \\
\texttt{A\_correcciones\_gateway\_20261110-1713.txt} & Correcciones al gateway: cliente HTTP, limpiezas, \texttt{/salud}; pruebas de B y de A4 después (2.1 y 6.3). \\
\texttt{A6.2\_simulador\_gateway\_20261110-1501.txt} (\texttt{.csv}) & 79 de 139 \texttt{/tareas/siguiente} con 503 por locks viejos (6.3). \\
\texttt{A6.2\_simulador\_gateway\_20261110-1721.txt} (\texttt{.csv}) & Esperas de 10~s por \texttt{database is locked} en SQLite y la medición de WAL (6.3). \\
\texttt{A6.2\_simulador\_gateway\_20261110-1726.txt} (\texttt{.csv}) & El mismo simulador sin errores después de las correcciones (6.3). \\
\texttt{A6.2\_simulador\_borradores\_20261114-2318.txt} (\texttt{.csv}) & El simulador con los borradores y latidos de B4.2 ya integrados: revisión en 200 y sin errores (6.4). \\
\texttt{P6\_carga\_20261110-1800.md} y su carpeta (\texttt{resumen.txt}, \texttt{desglose.txt}, \texttt{lago.txt}, CSV por corrida) & P6 final: latencias, envíos por minuto, \texttt{docker stats}, eventos en el lago y su atraso (6.2 y 6.3). \\
\texttt{P6\_carga\_20261110-1727.md} y su carpeta (\texttt{eventos.txt}) & Primera P6: el publicador de eventos atrasado (6.3). \\
\texttt{P6\_gateway\_2procesos\_20261110-1820/} & Experimento con el gateway en 2 procesos (6.3). \\
\texttt{P8\_mongo\_gateway\_20261110-1715.md} y sus \texttt{.txt}/\texttt{.csv} por escenario & P8 de MongoDB: RTO y RPO de los tres escenarios (7.1). \\
\texttt{P8\_mongo\_gateway\_20261110-1453.md}, \texttt{P8\_mongo\_20261022-2028.md} & P8 antes de corregir el gateway y directa al registro (7.1). \\
\bottomrule
\end{tabular}
\end{center}

\falta{B y C: agregar sus archivos (P1, P3, P8 de Valkey, P4, P5) a esta tabla}


\end{document}
